% ============================================================
%  CHAPTER 1 — INTRODUCTION
% ============================================================
\chapter{Introduction}

\section{Project Overview}

The rapid expansion of the Internet of Things (IoT) ecosystem has introduced an unprecedented scale of networked devices, generating billions of data packets daily across heterogeneous communication channels. Traditional network simulation environments model these interactions using static heuristics and predefined routing tables, fundamentally limiting their capacity to adapt to real-time anomalies, congestion events, and adversarial intrusions.

\textbf{Aegis-IoT} is a Neural Network Intelligence Platform that obliterates this paradigm. Built on the AnyLogic discrete-event simulation framework, this platform fuses natively transpiled machine learning kernels with real-time reinforcement learning and federated distributed training to create a self-healing, predictive network topography. The system identifies threats, optimizes traffic loads, and forecasts congestion—all in microseconds, not seconds.

\section{Problem Statement}

Modern IoT networks face three critical challenges that existing simulation platforms fail to address:

\begin{enumerate}[leftmargin=2cm]
    \item \textbf{Latency Prediction at Scale:} Current IoT simulators lack embedded predictive analytics. Network operators cannot anticipate congestion events before they manifest, resulting in reactive rather than proactive traffic management.
    
    \item \textbf{Real-Time Anomaly Detection:} Distributed Denial-of-Service (DDoS) attacks targeting IoT infrastructure exploit the resource-constrained nature of edge devices. Existing defenses rely on centralized, post-hoc analysis rather than inline, microsecond-scale threat neutralization.
    
    \item \textbf{Privacy-Preserving Intelligence:} Centralizing raw IoT telemetry data for model training violates data sovereignty principles and introduces single points of failure. A decentralized learning paradigm is essential.
\end{enumerate}

\section{Objectives}

The primary objectives of this project are:

\begin{itemize}[leftmargin=2cm]
    \item Design and implement a high-fidelity IoT network simulation using AnyLogic's agent-based discrete-event engine with MQTT protocol modeling.
    \item Train and deploy machine learning models (Random Forest, OneClassSVM, Autoregressive Forecaster) for real-time latency prediction, anomaly detection, and time-series forecasting.
    \item Eliminate Python runtime dependencies by transpiling trained models into native Java bytecode using the \texttt{m2cgen} framework for zero-latency inference.
    \item Implement a Tabular Q-Learning reinforcement learning agent for autonomous, adaptive load balancing within the simulation loop.
    \item Develop a Federated Learning pipeline implementing the FedAvg algorithm for privacy-preserving, distributed model training across simulated edge nodes.
    \item Simulate adversarial GAN-based DDoS traffic to stress-test the anomaly detection pipeline.
    \item Build a decoupled, four-monitor Network Operations Center (NOC) using hardware-accelerated Java Swing for real-time telemetry visualization.
    \item Develop a modern React/Vite web dashboard for cross-platform analytics and network topography mapping.
\end{itemize}

\section{Scope and Contributions}

This project makes the following technical contributions:

\begin{table}[H]
\centering
\caption{Summary of Technical Contributions}
\label{tab:contributions}
\begin{tabularx}{\textwidth}{@{}lX@{}}
\toprule
\textbf{Contribution} & \textbf{Description} \\
\midrule
AST Transpilation & Compilation of Python ML models into native Java decision trees via \texttt{m2cgen}, achieving sub-millisecond inference inside the AnyLogic JVM. \\
\addlinespace
RL Load Balancing & Implementation of a Tabular Q-Learning agent (\texttt{QTableBalancer.java}) with discrete state-action spaces for autonomous congestion management. \\
\addlinespace
Federated Learning & End-to-end FedAvg pipeline with SGD edge nodes and cloud-based weight aggregation, preserving data privacy. \\
\addlinespace
GAN DDoS Simulation & Adversarial traffic injection engine generating synthetic volumetric DDoS packets to validate the OneClassSVM radar. \\
\addlinespace
4-Window NOC & Decoupled Java Swing dashboards with hardware-accelerated Java2D rendering for real-time telemetry. \\
\addlinespace
Digital Twin & Composite virtual node health scoring combining battery depletion modeling and anomaly density tracking. \\
\addlinespace
Web Analytics & React/Vite frontend with Recharts and Framer Motion for cross-platform network visualization. \\
\bottomrule
\end{tabularx}
\end{table}

\section{Report Organization}

This report is organized as follows:

\begin{itemize}
    \item \textbf{Chapter 2} reviews the literature on IoT simulation, machine learning for network security, federated learning, and reinforcement learning.
    \item \textbf{Chapter 3} presents the system architecture, including the AnyLogic agent topology, data flow, and component interactions.
    \item \textbf{Chapter 4} details the machine learning pipeline: data preprocessing, model training, evaluation, and the \texttt{m2cgen} transpilation process.
    \item \textbf{Chapter 5} covers the Reinforcement Learning load balancer and the Federated Learning protocol.
    \item \textbf{Chapter 6} describes the four-window Java Swing NOC dashboard system.
    \item \textbf{Chapter 7} presents the React/Vite web dashboard architecture.
    \item \textbf{Chapter 8} explains the injection engine—the Python-based mutation scripts that dynamically patch AnyLogic project files.
    \item \textbf{Chapter 9} presents experimental results and performance analysis.
    \item \textbf{Chapter 10} concludes the report with a summary and future work.
\end{itemize}
